%!TEX root = main.tex

\section{Method}

% \subsection{Method overview}
Our goal is to learn visual features in an online fashion without supervision.
To that effect, we propose an online clustering-based self-supervised method.
Typical clustering-based methods~\cite{asano2019self,caron2018deep} are offline in the sense that they alternate between a cluster assignment step where image features of the entire dataset are clustered, and a training step where the cluster assignments, \ie, ``codes'' are predicted for different image views.
Unfortunately, these methods are not suitable for online learning as they require multiple passes over the dataset to compute the image features necessary for clustering.
In this section, we describe an alternative where we enforce consistency between codes from different augmentations of the same image.
This solution is inspired by contrastive instance learning~\cite{wu2018unsupervised} as we do not consider the codes as a target, but only enforce consistent mapping between views of the same image.
Our method can be interpreted as a way of contrasting between multiple image views by comparing their cluster assignments instead of their features.

More precisely, we compute a code from an augmented version of the image and predict this code from other augmented versions of the same image.
Given two image features $\mathbf{z}_t$ and $\mathbf{z}_s$ from two different augmentations of the same image, we compute their codes $\mathbf{q}_t$ and $\mathbf{q}_s$ by matching these features to a set of $K$ prototypes $\{\mathbf{c}_1,\dots, \mathbf{c}_K\}$.
We then setup a ``swapped'' prediction problem with the following loss function:
\begin{eqnarray}
L(\mathbf{z}_t, \mathbf{z}_s) & = &\ell(\mathbf{z}_t, \mathbf{q}_s) + \ell(\mathbf{z}_s, \mathbf{q}_t),
\label{eq:twoviews}
\end{eqnarray}
where the function~$\ell({\mathbf z},{\mathbf q})$ measures the fit between features ${\mathbf z}$ and a code ${\mathbf q}$, as detailed later.  
Intuitively, our method compares the features $\mathbf{z}_t$ and $\mathbf{z}_s$ using the intermediate codes $\mathbf{q}_t$ and $\mathbf{q}_s$. 
If these two features capture the same information, it should be possible to predict the code from the other feature.
A similar comparison appears in contrastive learning where features are compared directly~\cite{wu2018unsupervised}. 
In~\cref{fig:oto}, we illustrate the relation between contrastive learning and our method.

% More formally, the loss function between two transformations of the same image is the sum of two multinomial logistic loss function with a temperature parameter $\tau$:
% Basically, this loss discriminates the cluster assignment, or code of one transformation with the feature of the other transformation.
% This ensures that two different views of the same images have a consistent mapping to the code vectors.

\begin{figure}[t]
\centering
\begin{tabular}{c | c}
\includegraphics[height=.2\linewidth]{simclr_simple.pdf} &
\includegraphics[height=.2\linewidth]{oto_simple.pdf} \\
\\
Contrastive instance learning & \OURSFULL (Ours)
\end{tabular}
   \caption{\textbf{Contrastive instance learning (left) \vs\ \OURS (right).}
% Both methods obtain features from two different data augmentations of the same image.
In contrastive learning methods applied to instance classification, the features from different transformations of the same images are compared directly to each other.
In \OURS, we first obtain ``codes'' by assigning features to prototype vectors.
We then solve a ``swapped'' prediction problem wherein the codes obtained from one data augmented view are predicted using the other view.
Thus, \OURS does not directly compare image features.
Prototype vectors are learned along with the ConvNet parameters by backpropragation. 
}\label{fig:oto}
\end{figure}

\subsection{Online clustering}
\label{sec:proto}

Each image $\mathbf{x}_n$ is transformed into an augmented view $\mathbf{x}_{nt}$ by applying a transformation $t$ sampled from the set $\mathcal{T}$ of image transformations.
The augmented view is mapped to a vector representation by applying a non-linear mapping $f_\theta$ to $\mathbf{x}_{nt}$.
The feature is then projected to the unit sphere,~\ie,~ $\mathbf{z}_{nt} = f_\theta(\mathbf{x}_{nt}) / \|f_\theta(\mathbf{x}_{nt})\|_2$.
We then compute a code $\mathbf{q}_{nt}$ from this feature by mapping $\mathbf{z}_{nt}$ to a set of $K$ trainable prototypes vectors,  $\{\mathbf{c}_1,\dots,\mathbf{c}_K\}$.
We denote by $\mathbf{C}$ the matrix whose columns are the $\mathbf{c}_1, \dots, \mathbf{c}_k$. 
We now describe how to compute these codes and update the prototypes online.

\paragraph{Swapped prediction problem.} 
The loss function in Eq.~(\ref{eq:twoviews}) has two terms that setup the ``swapped'' prediction problem of predicting the code $\mathbf{q}_t$ from the feature $\mathbf{z}_s$, and $\mathbf{q}_s$ from $\mathbf{z}_t$. 
Each term represents the cross entropy loss between the code and the probability obtained by taking a softmax of the dot products of $\mathbf{z}_i$ and all prototypes in $\mathbf{C}$, \ie, 
\begin{equation}
  \ell(\mathbf{z}_t, \mathbf{q}_s) = - \sum_{k} \mathbf{q}_s^{(k)} \log \mathbf{p}_t^{(k)}, \quad \text{where} \quad \mathbf{p}_t^{(k)} = \frac{ \exp \left ( \frac{1}{\tau} \mathbf{z}_t^\top \mathbf{c}_k \right ) }{\sum_{k'} \exp \left ( \frac{1}{\tau} \mathbf{z}_t^\top \mathbf{c}_{k'} \right ) }.
  \label{eq:loss}
\end{equation}
where $\tau$ is a temperature parameter~\cite{wu2018unsupervised}.
Taking this loss over all the images and pairs of data augmentations leads to the following loss function for the swapped prediction problem:
\begin{equation}
  - \frac{1}{N} \sum_{n=1}^N \sum_{s,t \sim \mathcal{T}} \left [ 
  \frac{1}{\tau} \mathbf{z}_{nt}^\top \mathbf{C} \mathbf{q}_{ns}
  + \frac{1}{\tau} \mathbf{z}_{ns}^\top \mathbf{C} \mathbf{q}_{nt}
  - \log \sum_{k=1}^K \exp \left ( \frac{\mathbf{z}_{nt}^\top \mathbf{c}_k}{\tau} \right )  
  - \log \sum_{k=1}^K \exp \left ( \frac{\mathbf{z}_{ns}^\top \mathbf{c}_k}{\tau} \right ) \right ] .
  \nonumber
\label{eq:swapploss2}
\end{equation}
This loss function is jointly minimized with respect to the prototypes $\mathbf{C}$ and the parameters $\theta$ of the image encoder $f_\theta$ used to produce the features $(\mathbf{z}_{nt})_{n,t}$.

\paragraph{Computing codes online.}

In order to make our method online, we compute the codes using only the image features within a batch.
Intuitively, as the prototypes $\mathbf{C}$ are used across different batches, \OURS clusters multiple instances to the prototypes.
We compute codes using the prototypes $\mathbf{C}$ such that all the examples in a batch are equally partitioned by the prototypes. 
This equipartition constraint ensures that the codes for different images in a batch are distinct, thus preventing the trivial solution where every image has the same code.
Given $B$ feature vectors $\mathbf{Z}=[\mathbf{z}_{1},\dots,\mathbf{z}_{B}]$, we are interested in mapping them to the prototypes $\mathbf{C}=[\mathbf{c}_1,\dots,\mathbf{c}_K]$.
We denote this mapping or codes by $\mathbf{Q} = [\mathbf{q}_1, \dots, \mathbf{q}_B]$, and optimize $\mathbf{Q}$ to maximize the similarity between the features and the prototypes , \ie, 
\begin{equation}
  \label{eq:assign}
  \max_{\mathbf{Q}\in\mathcal{Q}} \ \text{Tr} \left( \mathbf{Q}^\top \mathbf{C}^\top \mathbf{Z} \right) +  \varepsilon H(\mathbf{Q}),
\end{equation}
where $H$ is the entropy function, $H(\mathbf{Q}) = -\sum_{ij} \mathbf{Q}_{ij} \log \mathbf{Q}_{ij}$ and $\varepsilon$ is a parameter that controls the smoothness of the mapping.
We observe that a strong entropy regularization (\ie~using a high $\varepsilon$) generally leads to a trivial solution where all samples collapse into an unique representation and are all assigned uniformely to all prototypes.
Hence, in practice we keep $\varepsilon$ low.
Asano \etal~\cite{asano2019self}  enforce an equal partition by constraining the matrix $\mathbf{Q}$ to belong to the transportation polytope.
They work on the full dataset, and we propose to adapt their solution to work on minibatches by restricting the transportation polytope to the minibatch:
\begin{equation}
  \mathcal{Q} = \left \{\mathbf{Q}\in\mathbb{R}_+^{K\times B} ~|~\mathbf{Q} \mathbf{1}_B = \frac{1}{K} \mathbf{1}_K, \mathbf{Q}^\top \mathbf{1}_K = \frac{1}{B} \mathbf{1}_B \right \},
\end{equation}
where $\mathbf{1}_K$ denotes the vector of ones in dimension $K$.
These constraints enforce that on average each prototype is selected at least $\frac{B}{K}$ times in the batch.

Once a continuous solution $\mathbf{Q}^*$ to Prob.~(\ref{eq:assign}) is found, a discrete code can be obtained by using a rounding procedure~\cite{asano2019self}.
%This code can be derived by resorting to hard inference or sampling treating each $\mathbf{q}_n^*$ as a probability.
Empirically, we found that discrete codes work well when computing codes in an offline manner on the full dataset as in Asano~\etal~\cite{asano2019self}.
However, in the online setting where we use only minibatches, using the discrete codes performs worse than using the continuous codes.
An explanation is that the rounding needed to obtain discrete codes is a more aggressive optimization step than gradient updates.
While it makes the model converge rapidly, it leads to a worse solution.
We thus preserve the soft code $\mathbf{Q}^*$ instead of rounding it. 
These soft codes $\mathbf{Q}^*$ are the solution of Prob.~(\ref{eq:assign}) over the set $\mathcal{Q}$ and takes the form of a normalized exponential matrix~\cite{cuturi2013sinkhorn}:
\begin{equation}
\label{eq:qstar}
  \mathbf{Q}^*= \text{Diag}(\mathbf{u}) \exp\left(\frac{\mathbf{C}^\top \mathbf{Z}}{\varepsilon} \right) \text{Diag}(\mathbf{v}),
\end{equation}
where $\mathbf{u}$ and $\mathbf{v}$ are renormalization vectors in $\mathbb{R}^K$ and $\mathbb{R}^B$ respectively.
The renormalization vectors are computed using a small number of matrix multiplications using the iterative Sinkhorn-Knopp algorithm~\cite{cuturi2013sinkhorn}.
In practice, we observe that using only $3$ iterations is fast and sufficient to obtain good performance.
%We conjure that it is the same reason that rounding codes to discrete values deteriorate the quality of our model by converging too rapidly. 
Indeed, this algorithm can be efficiently implemented on GPU, and the alignment of $4$K features to $3$K codes takes $35$ms in our experiments, see~\cref{sec:exp}.
%Interestingly, the parameter $\varepsilon$ plays a similar role of controlling the smoothness of the exponential in Eq.~(\ref{eq:qstar}) as the temperature $\tau$ in Eq.~(\ref{eq:loss}).

\paragraph{Working with small batches.}
When the number $B$ of batch features is too small compared to the number of prototypes $K$, it is impossible to equally partition the batch into the $K$ prototypes.
% The mapping quality of the batch features $\mathbf{Z}$ to the codes $\mathbf{C}$ becomes poor .
Therefore, when working with small batches, we use features from the previous batches to augment the size of $\mathbf{Z}$ in Prob.~(\ref{eq:assign}).
%This is similar to the memory bank mechanism used in contrastive methods~\cite{he2019momentum,wu2018unsupervised}.
Then, we only use the codes of the batch features in our training loss.
In practice, we store around $3$K features, \ie, in the same range as the number of code vectors.
This means that we only keep features from the last $15$ batches with a batch size of $256$, while contrastive methods typically need to store the last $65$K instances obtained from the last $250$ batches~\cite{he2019momentum}.

%\paragraph{Relation to contrastive learning.}
%Contrastive learning~\cite{hadsell2006dimensionality} compares the two features $\mathbf{z}_t$ and $\mathbf{z}_s$, and contrasts this similarity against features from other images.
%This is an extreme case of our formulation where the centroids are previously-stored image features and the code is a direct mapping between the current image and its stored representation~\cite{wu2018unsupervised}.
%\jm{The previous sentence needs a small clarification as it is a dangerous one. We need to emphasize that this extreme case is not what we do: centroids are not previously-stored image features, and we have good reasons not doing so. }
%As the number of classes grows with the dataset, it is standard to approximate the loss with a subset of randomly selected centroids, or images~\cite{chen2020simple, he2019momentum, wu2018unsupervised}.
%\jm{Again, this may give the impression that procedure applies to our method.}
%In particular, contrastive methods like MemoryBank~\cite{wu2018unsupervised}, MoCo~\cite{he2019momentum} are designed to work with small batches.
%MemoryBank requires to store features for the entire dataset and MoCo needs extra memory for a Momentum Encoder and for storing around $65$K instances which means features that are $250$ batches away.
%\jm{maybe clarify what it means to be 250 batches away.}
%In contrast, our method does not require an extra encoder and only needs to store features in the same range as the number of code vectors, which is typically around $3$K (therefore only $15$ batches away).
\begin{figure}[t!]
  %\setlength{\tabcolsep}{3.5pt}
\begin{minipage}{0.45\linewidth}
\centering
    \begin{tabular}{@{} l l c c @{}}
      \toprule
      Method      & Arch. & Param. & Top1  \\
      \midrule
      Supervised  & \resnetfifty    & 24  & 76.5  \\
%      		  & EffNet-B7    & 66  & 84.4  \\
      \midrule
      Colorization~\cite{zhang2016colorful}     & \resnetfifty  & 24 & 39.6 \\
      Jigsaw~\cite{noroozi2016unsupervised}     & \resnetfifty  & 24 & 45.7 \\
      NPID~\cite{wu2018unsupervised}   & \resnetfifty  & 24 & 54.0   \\
      BigBiGAN~\cite{donahue2019large} & \resnetfifty  & 24 & 56.6   \\
      LA~\cite{zhuang2019local}        & \resnetfifty  & 24 & 58.8 \\
      NPID++~\cite{misra2019self}      & \resnetfifty  & 24 & 59.0   \\
      MoCo~\cite{he2019momentum}       & \resnetfifty  & 24 & 60.6   \\
      SeLa~\cite{asano2019self}        & \resnetfifty  & 24 & 61.5   \\
      PIRL~\cite{misra2019self}        & \resnetfifty  & 24 & 63.6   \\
      CPC v2~\cite{henaff2019data}     & \resnetfifty  & 24 & 63.8 \\
      PCL~\cite{li2020prototypical}     & \resnetfifty  & 24 & 65.9 \\
      SimCLR~\cite{chen2020simple}     & \resnetfifty  & 24 & 70.0  \\
      MoCov2~\cite{chen2020improved}   & \resnetfifty  & 24 & 71.1  \\
\midrule
      \OURS   		       & \resnetfifty  & 24 & \bf{75.3}        \\ 
      \bottomrule
 \end{tabular}
\end{minipage}
~~~~~\begin{minipage}{0.5\linewidth}
\centering
\includegraphics{figure_param_with_w5.pdf} 
\end{minipage}
    \caption{
      \textbf{Linear classification on \ImNet.} 
Top-1 accuracy for linear models trained on frozen features from different self-supervised methods.
\textbf{(left)}
Performance with a standard ResNet-50.
\textbf{(right)} Performance as we multiply the width of a ResNet-50 by a factor $\times2$, $\times4$, and $\times5$.
}
    \label{fig:linear_inet}
\end{figure}



\subsection{Multi-crop: Augmenting views with smaller images}
\label{sec:multicrop}
As noted in prior works~\cite{chen2020simple,misra2019self}, comparing random crops of an image plays a central role by capturing information in terms of relations between parts of a scene or an object. 
Unfortunately, increasing the number of crops or ``views'' quadratically increases the memory and compute requirements. 
We propose a \texttt{multi-crop} strategy where we use two standard resolution crops and sample $V$ additional low resolution crops that cover only small parts of the image. 
Using low resolution images ensures only a small increase in the compute cost. 
% Additionally, we observe that mapping these small parts of an image to more global views increases performance.
Specifically, we generalize the loss of Eq~(\ref{eq:twoviews}):
\begin{equation}
L(\mathbf{z}_{t_1}, \mathbf{z}_{t_2}, \dots, \mathbf{z}_{t_{V + 2}}) = \sum_{i \in \{1, 2\}} \sum_{v=1}^{V + 2} \mathbf{1}_{v \neq i} \ell(\mathbf{z}_{t_v}, \mathbf{q}_{t_i}).
\end{equation}
Note that we compute codes using only the full resolution crops.
Indeed, computing codes for all crops increases the computational time and we observe in practice that it also alters the transfer performance of the resulting network.
An explanation is that using only partial information (small crops cover only small area of images) degrades the assignment quality.
Figure~\ref{fig:multicrop-exp} shows that \texttt{multi-crop} improves the performance of several self-supervised methods and is a promising augmentation strategy.
%Downsizing the resolutions of the crops does not affect transfer performance while greatly reducing the compute requirement.
%Training with $160 \times 160$ images even boosts the performance by $+0.5\%$ on ImageNet while halving the compute and memory costs compared to the standard $224 \times 224$ images.

